\section{Design}

\aebench is a framework for running and evaluating AI coding agents on artifact evaluation tasks. 
The main feature of \aebench is running the agents and recording runtime information
  to allow to validate whether or not the artifact actually ran the artifact and 
  didn't fake the results by copying them from the paper or artifact repository 
  through repository history.
  
\aebench runs an agent in a controlled environment, allows it to attempt the documented 
  reproduction procedure documented in the README file of the artifact, evaluates the result 
  using case specific checks which we called oracles, and checks the commands that 
  the agent runs against the README.

We design \aebench to implement both the execution framework and the evaluation framework
  but the two are separated and agents cannot see the evaluation stage.

\begin{itemize}
  \item \aebench execution framework prepares the artifact, launches an AI agent, manages its environment, and records what happens.
  \item \aebench evaluation framework independently checks whether the agent successfully reproduced the artifact's expected behavior or results.
\end{itemize}


At a high level, AEBench has seven functional components.

\begin{enumerate}
  \item CLI and configuration component for user commands, settings, case selection, agent/model options

  \item Artifact loading and validation which includes artifact configuration, artifact source, README, reproduction criteria

  \item Artifact running / runtime and orchestration which coordinates preparation, agent execution, oracle evaluation, and results saving

  \item Runtime backend which has the local or Docker execution, workspace and container lifecycle

  \item Agent running logic, including prompts, credentials, agent commands, and configuration for specific agent

  \item Oracle evaluation that check environment, build, preparation, and experiment results

  \item Results, logs, and traces which all go into a llm-as-a-judge module that checks whether the agent has actually 
\end{enumerate}